\section{Learning from extremal prices}\label{sec:prices}

In the LP algorithm, a query chooses scales for the coordinates and returns
signs at those scales. The coordinates where a nonnegative circuit has its largest
scaled entries receive the sign of its side. Theorem~\ref{thm:prices} turns this limited
feedback into one price vector satisfying every consistent circuit
hypothesis. We prove it in the abstract setting, with no access to the
subspaces that generated the signs.

For integers \(N,R\ge1\), a hypothesis specifies integer demands
\(a_s\in[-R,R]\) on a nonempty set \(S\) of \emph{slots}, the signed
coordinates in \(\Slots=[N]\times\{-1,+1\}\). For a vector \(y\)
write \(y_s=h y_i\) at slot \(s=(i,h)\). The surplus
\(d_s(y)=a_s-y_s\) is the amount by which the price falls short of
the demand. A query at \(y\) returns labels \(z\in\{\pm1\}^N\).
Consistency requires \(z_s=1\) at \emph{every} slot attaining
\(\max_{u\in S}d_u(y)\) whenever this maximum is nonnegative.
The desired output has \(d_s(y)\le0\) on all slots of every hypothesis
consistent with the entire sequence of answers.

Consider a hypothesis with two demands on positive slots. If the prices
fall short by five and two units, consistency requires a positive label
on the first slot and permits either sign on the second. An expert
guessing only the second demand could act on the wrong sign. Suppose
that violations carry enough of a chosen probability distribution on
the slots. A successful recursive call will then give correct feedback
on a set of fixed positive mass. An expert sampling four slots from
that distribution has a good chance to see one of them.

\subsection{A market for feedback}

In one round, consider a finite pool of experts \(e\),
each with wealth \(w_e>0\).
An expert comes with a finite list of options. An option has utility
\(b-y_s\) at prices \(y\), for some bid \(b\) and slot \(s\);
each expert also has an outside option with utility zero.
Denote the maximum utility over expert \(e\)'s options by \(U_e(y)\).

Choose prices \(y\in\RR^N\) minimizing
\[
        \Phi(y)=\sum_e w_e U_e(y) .
\]
Introducing an upper bound on each expert's option utilities expresses
this as a feasible linear program. The outside options bound its
objective below by zero, so it attains a minimum. With rational data
we may take a rational minimizer. When the feedback is
\(v\in\{\pm1\}^N\), pay each maximizing slot option \(v_s\), and pay
the outside option zero. Define \(g_e\) to be the \emph{least} payoff
among expert \(e\)'s maximizing options. The one-sided directional
derivative of its utility satisfies \(U'_e(y;v)=-g_e\), since the derivative of a maximum
is the largest active slope in the given direction. Minimality of
\(\Phi\) gives, simultaneously for every feedback vector,
\begin{equation}\label{eq:market}
         \sum_e w_e g_e=-\Phi'(y;v)\le0 .
\end{equation}
Consequently the update \(w_e\gets w_e(1+\eta g_e)\), with rational
\(0<\eta<1\), preserves positivity and never increases total wealth.

There is a market interpretation of~\eqref{eq:market}. At a minimizer,
zero is a subgradient of \(\Phi\). Each expert's wealth can therefore
be divided among its maximizing options so that the two signs of each
coordinate receive equal total weight. This division would give total
gain zero for every \(v\). Taking the least payoff at each tie gives
the inequality directly. Compare the market-clearing condition in
matching markets~\cite[Chapter~10]{EK} and the halfspace condition
in Blackwell approachability~\cite{Black}. For multiplicative weights
and boosting, see~\cite{Litt,FS,AHK}.

\paragraph{One expert per hypothesis.}
Give each hypothesis initial wealth one and bids \(a_s\) on its slots.
In each round query labels at the minimizing prices and use them as
feedback in the update. A hypothesis consistent with all answers has
gain one whenever any of its demands is violated. Indeed, the
maximum surplus is then positive and every maximizing option is a
correctly labeled slot. In every other round its gain is nonnegative:
a maximizing slot tied with the outside option has surplus zero
and is also correctly labeled. Thus \(k\) rounds with an unmet demand
give this expert wealth at least \((1+\eta)^k\). The initial total
wealth \(M\), the number of hypotheses, bounds that quantity.

Run an odd number of rounds greater than \(2\log_{1+\eta}M\).
Each consistent hypothesis has a majority of rounds meeting all of
its demands, although the majority may differ among hypotheses.
Taking coordinatewise medians meets them all because a signed
coordinate of the median is the median of that signed coordinate.
There are at most \((2R+2)^{2N}\) hypotheses, so for fixed \(\eta\)
the query count is \(O(N\log(R+1))\).

For demands \(a_s=p_s\) on both signs of every coordinate, the
expert's utility is \(\|p-y\|_\infty\). The minimizing prices are
then weighted geometric medians in this norm. Chen, Li and
Yannakakis~\cite{CLY} use related balanced points to bound query
complexity for contractions (see also~\cite{HLSW,FH}). At \(y=p\)
both signs tie at zero, so these exact point demands cannot be
consistent with the answer there. Section~\ref{sec:contractions}
uses approximate demands in applying the theorem.

Explicitly maintaining an expert for every hypothesis costs exponential
time. With four guessed slots per expert, the implementation in
Section~\ref{sec:implement-pr} will need only a constant-dimensional
space of guesses. The remaining question is how these experts can
profit if they miss the one slot at which a label is guaranteed.

\subsection{Amplifying the feedback}\label{sec:amplify}

An internal node makes an odd number of \emph{probes}. Each probe
compares two price vectors, its own proposal and a price vector found
by a recursive call.

\begin{enumerate}
 \item Choose proposal prices \(y\) (ultimately this will be by a market).
 Call a fresh child with \(y\) added to the price history,
 the list of proposals along the path from the root.
 \item Let \(x\) be the prices returned by that child.
 Query labels \(z\) at \(x\), and form feedback at \(y\):
 \begin{equation}\label{eq:feedback}
 v_i =
 \begin{cases}
     \sgn(x_i-y_i) & x_i\ne y_i,\\
     z_i & x_i=y_i.
 \end{cases}
 \end{equation}
\end{enumerate}
After all probes, return the coordinatewise median of the prices
returned by the children. We will see how to end the recursion
after logarithmically many levels.

\paragraph{Assigning thresholds.}
Fix the complete run and a hypothesis \((S,a)\) consistent with it.
Assign each node a threshold \(t\ge0\). A node \emph{succeeds} for this
hypothesis if its returned prices make every surplus at most \(t\).
The root's threshold is zero. The algorithm uses only the price history
and the queried labels. We may choose these thresholds separately in
the proof for each consistent hypothesis.

Suppose a child returns \(x\) and succeeds at its own threshold \(t'\).
Then the feedback at its parent's proposal \(y\) obeys

\begin{equation}\label{eq:correct}
            d_s(y)\ge t' \ \Longrightarrow\ v_s=1
            \qquad(s\in S).
\end{equation}
Indeed, \(d_s(x)\le t'\le d_s(y)\) gives \(x_s\ge y_s\).
Strict inequality already gives \(v_s=1\);
in case of equality the slot must maximize surplus at \(x\)
with value \(t'\),
so that \(z_s=1\) by consistency.

Thus if many surpluses at \(y\) exceed \(t\), we can
raise the child's threshold yet still aim to get feedback
agreeing with many slots.
Here is how we do the assignment, propagating also a probability
distribution \(m\) on \(S\), uniform at the root.
For a node with \(m,t\) and a probe there, represent \(m\) by subdividing \([0,1)\),
with a piece of length \(m_s\) for each slot in order of
nonincreasing surplus \(d_s(y)\) (break ties deterministically).
Call the probe
\begin{itemize}
 \item \emph{light} if \(d_s(y)>t\)
 on mass at most \(\delta=1/5\). Then give the child the threshold \(t'=t\).
 \item \emph{heavy} otherwise. Give the child as threshold
 \(t'\) a high surplus quantile strictly above \(t\):
 as \emph{pivot} take any slot meeting the band
 \[
           [\gamma,\delta),\qquad \gamma=3/16
 \]
 in positive length, and let \(t'\) be its surplus.
 (Later we will make a specific choice in this band.)
 Thus
 \begin{equation}\label{eq:band}
 m\{s:d_s(y)\ge t'\}\ \ge\ \gamma,
 \qquad
 m\{s:d_s(y)> t'\}\ \le\ \delta .
 \end{equation}
 See Figure~\ref{fig:band}.
\end{itemize}

A threshold is recorded either as zero or by the pivot slot and the
level of the proposal at which its surplus gave the threshold.
A light probe passes this description to the child. A heavy probe
replaces it by the pivot and the current level.

To form the child's distribution \(m'\), multiply the first
\(\delta\) of the sorted mass by four and divide the resulting mass
by \(1+3\delta\). A slot crossing the cutoff receives a partial
boost. Every slot with \(d_s(y)>t'\) receives the full boost, by
the light-probe rule or by~\eqref{eq:band}. Once \(m\) and \(y\)
are fixed, the same \(m'\) works for every allowed pivot. As with
the thresholds, the distributions are used only in the proof.

\begin{figure}[htbp]
\centering
\begin{tikzpicture}[
  x=35cm,y=.86cm,font=\small,
  slot/.style={black,line width=.87pt},
  guide/.style={black!48,densely dashed,line width=.4pt},
  axis/.style={black!55,line width=.45pt},
  bracket/.style={line width=.8pt}]
 % Show the first .35 of the unit mass interval at its actual scale.
 \fill[black!13] (.1875,0) rectangle (.20,3.38);
 \draw[guide] (0,1.85) -- (.35,1.85) node[right,text=black] {\(t'\)};
 \draw[guide] (0,1.12) -- (.35,1.12) node[right,text=black] {\(t\)};
 \node[anchor=east] at (0,3.05) {\(d_s(y)\)};
 \foreach \a/\b/\h in {
   0/.04/3.05,.04/.095/2.75,.095/.15/2.42,
   .15/.19375/1.85,.19375/.25/1.85,.25/.285/1.46,
   .285/.325/.78,.325/.35/.50}
   \draw[slot] (\a,\h) -- (\b,\h);
 % A vertical tick separates the two slots with the same surplus.
 \draw[slot] (.19375,1.74) -- (.19375,1.96);
 \fill (.1895,1.85) circle[radius=1.5pt];

 \draw[axis] (0,0) -- (.35,0);
 \foreach \x/\lab in {0/0,.10/0.10,.1875/\gamma,.20/\delta,.30/0.30,.35/0.35}{
   \draw[axis] (\x,0) -- (\x,-.1);
   \node[below=2pt,inner sep=1pt] at (\x,0) {\(\lab\)};
 }
 \node[anchor=north,font=\footnotesize] at (.296,-.44) {Cumulative slot mass};
 \draw[bracket,lagoon] (0,3.52) -- (0,3.68) -- (.25,3.68) -- (.25,3.52);
 \node[above=1pt,text=lagoon] at (.125,3.68) {\(v_s=1\) if the child succeeds};
 \draw[bracket,ember] (0,-.52) -- (0,-.68) -- (.20,-.68) -- (.20,-.52);
 \node[below=1pt,text=ember] at (.10,-.68) {Multiply this mass by four};
\end{tikzpicture}
\caption{A heavy probe, with the first \(0.35\) of the slot mass shown.
Each horizontal segment represents a slot, ordered by surplus. The dot selects
a pivot in the shaded band \([\gamma,\delta)=[3/16,1/5)\).
Two adjacent slots have the pivot surplus \(t'\). Success at \(t'\) guarantees
correct feedback on mass \(0.25\), including both tied slots.
Reweighting multiplies the first \(0.20\) of the mass by four before normalization.}
\label{fig:band}
\end{figure}


\paragraph{Leaves.}
Take the root at level zero, all leaves at level \(D\),
the least positive odd number with
\[
               (5/4)^D>2N .
\]
At a leaf, return the coordinatewise median of the \(D\)
price vectors in the history.

\begin{lemma}\label{lem:leaf}
Each leaf succeeds: its returned prices achieve maximum surplus
at most its threshold, for any assignment as above.
\end{lemma}
\begin{proof}
Otherwise surplus at some slot \(s\) exceeds the leaf's threshold
at the median. This also happened at more than half of the
\(D\) proposals on the path.
Since thresholds never decrease, at each such proposal \(y\)
we had \(d_s(y)>t'\), with \(t'\)
the threshold assigned to the child. Thus the slot's weight was boosted by
a factor of at least \(4^{D/2}\) over the path before accounting for
normalizations. Starting from mass \(1/|S|\),
the slot has ended with at least
\[
       \frac1{2N}\frac{4^{D/2}}{(1+3\delta)^D}>1 ,
\]
a contradiction.
\end{proof}

Recursive amplification with reweighting also appears in the
contraction algorithm of Soleymani and Farina~\cite[\S6]{SF} and the
mean-payoff game algorithms of~\cite{OAIDe,OAIRa,OAITu}. In
\cite[\S\S3--5]{OAIDe}, near-extreme deviations of hidden vectors imply
shifted labeling requirements deeper in the recursion. That work
also discusses the relation to Parys's recursion~\cite{Parys}.

\subsection{Betting on a few slots}

If a majority of probes are light and their children succeed, the
parent succeeds because those children met the parent's own threshold.
A successful child in a heavy probe gives a different guarantee:
feedback is correct on the slots at or above the pivot surplus, a set
of mass at least \(\gamma\). Four independent draws from \(m\) hit
that set with probability at least \(1-(1-\gamma)^4>1/2\).
This probability is the reason four slots suffice.

The algorithm uses the same pool \(E\) in every probe. An expert
specifies the following data.

\begin{itemize}
 \item four slots in \(\Slots\) (an ordered tuple),
 guessing an integer demand in \([-R,R]\) on each;
 \item either threshold zero, or a level \(j\in\{0,\ldots,D-1\}\)
 and pivot slot in \(\Slots\), again guessing such a demand on that slot.
\end{itemize}
Include all such experts (repetitions of slots allowed, even with
different demands).
Zero-threshold experts always bid; those who specify \(j\)
only bid at levels deeper than \(j\)
(otherwise they have only the outside option, giving zero gain).

In bidding on its four slots, an expert shifts
its guessed demands by \emph{subtracting} its threshold estimate \(\tau\).
For an expert specifying \(j\),
\(\tau\) is the surplus
on its pivot slot at the proposal at level \(j\) in the price history,
using its guess for the pivot demand; \(\tau\) may even be negative.
For zero-threshold experts \(\tau=0\).
Thus bidding gives the four option utilities
\[
                b-\tau-y_s
\]
(one for each tuple position, with slot \(s\) and guess \(b\)).
As before, experts keep the outside option of utility zero,
and gains use the feedback \(v\) with worst tie breaking.

We analyze again with a consistent hypothesis and thresholds
assigned as above.
At each node consider a distribution \(\nu\)
of experts: sample the four slots independently from \(m\),
match this node's threshold description (including pivot
and level if used), and use true demands throughout.
Thus for these experts \(\tau=t\)
and sampled slot utilities are \(d_s(y)-t\).

\begin{lemma}\label{lem:gain}
Fix a consistent hypothesis and any threshold assignment as in
Section~\ref{sec:amplify}, and set \(\eta=1/16\).
In any probe where the child succeeds, for market gains
at the proposal \(y\) we have
\[
 \sum_e \nu_e \ln(1+\eta g_e)\ge
 \begin{cases}
   0 & \text{if light},\\
   \epsilon:=1/256 & \text{if heavy}.
 \end{cases}
\]
Here \(\nu\) is from the parent node.
\end{lemma}
\begin{proof}
In the light case with \(t'=t\),
for these experts every maximizing option gets
nonnegative feedback by~\eqref{eq:correct}
(including zero on the outside option).

In the heavy case with \(t'>t\),
if the tuple hits \(\{s:d_s(y)\ge t'\}\),
the gain is one: every maximizing option has surplus at least
\(t'\) and hence correct sign, with no tie from the outside option.
This happens for mass at least \(1-(1-\gamma)^4\) of~\(\nu\).
Since gains are always at least \(-1\),
mean gain is at least
\(1-2(1-\gamma)^4>1/8\).
The bound follows from \(\ln(1+\eta g)\ge\eta g-\eta^2\)
for these gains.
\end{proof}

Consider running these markets independently at every internal node,
with unit initial wealth and an odd number of probes greater than
\((2/\epsilon)\ln|E|\). Suppose all the children succeed. If half or
more of the probes are heavy, then Lemma~\ref{lem:gain} makes the
final log wealth, averaged under the fixed parent distribution
\(\nu\), greater than \(\ln|E|\). Yet no expert can exceed the
initial total wealth \(|E|\). A majority of the probes must be
light, so the median rule gives success. Lemma~\ref{lem:leaf} starts
this induction at the leaves.

The resulting tree has logarithmic depth and \(O(\log |E|)\) children
at each internal node. To obtain a polynomial bound on its size we
reduce the branching to a constant below the root. A child may then
miss its threshold. The following wealth update lets its parent still
make progress on that probe.

\subsection{Investing recursively}

Give the child a small share of each expert's wealth and retain the
rest for the parent's market. If the child misses its threshold, we
will show that some experts earn sufficiently large returns there
to cover what they might lose on the feedback it supplies. Use the
tree of Section~\ref{sec:amplify}, the expert pool \(E\), the rate
\(\eta=1/16\), and absolute constants

\begin{equation}\label{eq:probepars}
      c=2^{-14},\qquad r=2^{19}+1 .
\end{equation}
Give each internal node \(r\) probes except the root, which gets
\[
       r_0=r+3072\lceil\log(|E|+1)\rceil .
\]

Each node receives as input the price history and
a vector giving each expert positive wealth.
At the root, start with wealth one each.
Each node returns prices as well as a new wealth vector.
A leaf simply returns median prices from the history,
and returns the input wealth unchanged.

At an internal node, use the probes successively, starting with
the input wealth and keeping any changes across probes.
In each probe, split the current wealth \(w\) as shown in
Figure~\ref{fig:invest}:
\begin{enumerate}
 \item Use \((1-c)w\) in the market:
  take proposal prices \(y\) minimizing \(\Phi\) with that wealth.
  \item Call the child with \(y\) added to the history
  and with wealth \(c w\);
  let \(x\) and \(w'\) be the prices and wealth returned by the child.
  \item Query labels at~\(x\) and take feedback~\eqref{eq:feedback}.
  Replace \(w_e\) for each expert~\(e\) by
  \begin{equation}\label{eq:invest}
       (1-c)w_e(1+\eta g_e) + w'_e ,
  \end{equation}
  using the gains in this probe's market.
 \end{enumerate}
Return median prices from the children,
and the wealth vector after the last replacement.

Output the prices returned by the root.
We leave implementation of the markets (including deterministic
choices of minimizers) for later in this section.

\begin{figure}[htbp]
\centering
\begin{tikzpicture}[
  x=1cm,y=1cm,font=\small,
  box/.style={draw=black!48,fill=white,rounded corners=2pt,align=center,
              inner xsep=6pt,inner ysep=5pt},
  cash/.style={draw=ember,line width=.88pt,rounded corners=3pt,
               -{Stealth[length=4.8pt]}},
  data/.style={draw=lagoon,dashed,line width=.72pt,-{Stealth[length=4.8pt]}},
  note/.style={font=\footnotesize,inner sep=2pt}]
 \coordinate (split) at (1.1,0);
 \node[box] (market) at (4.3,1.56) {Local market\\[-1pt]minimize \(\Phi\)};
 \node[box] (child) at (4.3,-1.56) {Child call};
 \node[box] (feedback) at (8.25,0) {Query labels at \(x\)\\[-1pt]form feedback by \eqref{eq:feedback}};
 \node[box] (gain) at (11.25,1.56) {Apply market gain\\[-1pt]\(\times\,(1+\eta g_e)\)};
 \node[circle,draw=ember,text=ember,line width=.65pt,inner sep=1.6pt]
       (sum) at (13.5,0) {\(+\)};

 \draw[cash] (0,0) -- (split) node[midway,above,note] {\(w_e\)};
 \fill[ember] (split) circle[radius=1.5pt];
 \draw[cash] (split) -- (1.1,1.56) -- (market.west);
 \node[above,note] at (2.05,1.56) {\((1-c)w_e\)};
 \draw[cash] (split) -- (1.1,-1.56) -- (child.west);
 \node[below,note] at (2.23,-1.56) {\(c w_e\)};
 \draw[cash] (market.east) -- (gain.west);
 \draw[cash] (gain.east) -- (13.5,1.56) -- (sum.north);
 \draw[cash] (child.east) -- (13.5,-1.56)
             node[midway,below,note] {Returned wealth \(w'_e\)} -- (sum.south);
 \draw[cash] (sum.east) -- (14.95,0)
             node[midway,above,note] {\(w_e^{\rm new}\)};

 \draw[data] (market.south) -- (child.north)
             node[midway,left=2pt,note] {\(y\)};
 \draw[data] (market.south east) -- (feedback.north west)
             node[midway,above right=0pt,note] {\(y\)};
 \draw[data] (child.north east) -- (feedback.south west)
             node[midway,below right=0pt,note] {\(x\)};
 \draw[data] (feedback.east) -- (gain.south)
             node[midway,below right=1pt,note] {\(v\)};
\end{tikzpicture}
\caption{Wealth and prices in one probe. Solid arrows track the two shares of
expert \(e\)'s wealth. Dashed arrows show how the proposal \(y\) and the
child's price vector \(x\) determine the feedback \(v\) for the local market.
The sum \(w_e^{\rm new}=(1-c)w_e(1+\eta g_e)+w'_e\) funds the next probe
or, after the last probe, returns to the parent.}
\label{fig:invest}
\end{figure}


Every node returns total wealth at most its input total. At a leaf the
two totals agree. At an internal node~\eqref{eq:market} bounds the
wealth paid by the market by its initial share, while the same claim
for the child bounds its returned share. Recombining the two bounds
and iterating over probes proves the claim.

\subsection{Profits cover failures}

Fix a consistent hypothesis. We will prove that a failed node gives a
large return on its input wealth, averaged over experts whose demands
agree with that hypothesis. This assertion is local: those experts may
have had little wealth when they entered the node. At the root, however,
every expert starts with wealth one and the total is at most \(|E|\).
There a sufficiently large return is impossible.

At each node let
\[
  \mu_e = \frac{\text{returned wealth for }e}
                     {\text{input wealth for }e},
       \qquad L_e = \ln(1+(c/2)\mu_e)
\]
for each expert.
Use \(L_e\) to count the bonus from returns here,
as a contribution to log wealth one level up.
Indeed, if in a probe the child has such values
\(\mu'_e,L'_e\), the wealth multiplier from~\eqref{eq:invest} is
\[
    (1-c)(1+\eta g_e)+c\mu'_e \ \ge\
    (1-c)(1+\eta g_e)\cdot (1+(c/2)\mu'_e)
\]
as \((1-c)(1+\eta g_e)<2\).
Consequently for any internal node,
\begin{equation}\label{eq:bonus}
 L_e\ \ge\ -\ln(2/c)+
     \sum_{\text{probes}}
     \big(\ln(1-c) + \ln(1+\eta g_e)+L'_e\big),
\end{equation}
by taking logs and summing to bound \(\ln\mu_e\).
For each term, the primed value is from that probe's child.
Each \(L'_e\) is nonnegative, even if the child's investment lost
wealth. Thus \(-\ln(1-c)\) bounds the loss per probe from sending
wealth to the child, while a large child return gives an additional
positive term.

To use a child's return in~\eqref{eq:bonus}, we must average its
bonuses with the parent's expert weights. The distribution \(\nu\)
from Lemma~\ref{lem:gain} samples four slots and specifies the
\emph{current} threshold. A heavy probe gives the child a new pivot
description, so its \(\nu'\) can be disjoint from the parent's \(\nu\).
We reserve weight at the parent for experts with each possible future
pivot level, including the one at which this probe takes place.

\paragraph{Comparing expert weights across one probe.}
For a fixed assignment of thresholds, give each node at level \(l\)
the following nonnegative weights \(\pi\) on the experts.

\begin{itemize}
 \item Include \(\nu\):
 the distribution with tuple slots i.i.d.\ drawn from~\(m\),
 matching this threshold description and using true demands.
 \item For each possible \emph{future} pivot level
 \(j\in\{l,\ldots,D-1\}\),
 include mass \(2^{l-j}\), distributed by picking four slots
 and the pivot slot at that level all independently from \(m\),
 with true demands. These experts have zero gains in the markets
 at this node, because they bid only deeper in the tree.
\end{itemize}
Thus~\(\pi\) is nonnegative of total mass at most three.

\begin{lemma}\label{lem:transfer-p}
For a complete run and a consistent hypothesis, there is an
assignment (choosing all thresholds and pivots as allowed in
Section~\ref{sec:amplify}) such that in every probe
\[
        \sum_e \pi'_e L'_e \ \le\ C \sum_e \pi_e L'_e ,
         \qquad C=80(5/2)^4=3125
\]
where \(\pi,\pi'\) refer to the parent and child respectively,
and \(L'\) is the vector of bonuses at the child.
\end{lemma}
\begin{proof}
Proceed down the tree. With \(m,t,\pi\) already determined for a
node, consider choosing the pivot in any heavy probe here at random:
take a uniform point in the band
\([\gamma,\delta)\) to determine the pivot slot.
This spreads out the choice: for any slot, the probability
is at most \(1/(\delta-\gamma)=80\) times that under \(m\).

We claim \(\EE\pi'\le C\pi\) componentwise, including a light probe
with its prescribed threshold. The child distribution \(m'\) is already
fixed before choosing the pivot. Since \(m'\le4m/(1+3\delta)=(5/2)m\),
sampling the four tuple slots costs a factor of at most \((5/2)^4\).

\begin{itemize}
 \item For the child's mass with its current threshold description:
if light, compare with the parent's such mass,
since the description does not change.
If heavy, compare instead with the parent's future mass for
\(j=l\), using the extra factor of 80 for the chosen pivot.
 \item For the child's future mass at any pivot level \(j\ge l+1\),
compare with the parent's future mass for that level.
Beyond the factor for the four slots there is a factor two
for total mass for that level, and at most \(5/2\)
for the pivot slot probabilities.
\end{itemize}
The parent's current description and its reserve for each future
level are distinct expert classes. The componentwise comparisons can
therefore be combined with the single factor \(C\).

The entire run was fixed before choosing pivots, so the child bonuses
\(L'_e\) are fixed nonnegative numbers. Multiplying the componentwise
expectation bound by \(L'_e\) and summing shows that at least one allowed
pivot satisfies the claimed inequality. Fix such a choice in every
heavy probe at this level and continue to the next level.

\end{proof}

\begin{lemma}\label{lem:charge}
For a consistent hypothesis, assign thresholds as in
Section~\ref{sec:amplify}, with the pivots chosen as in
Lemma~\ref{lem:transfer-p}.
If an internal node fails and has \(T\) probes, then
\begin{equation}\label{eq:charge}
           \sum_e \pi_e L_e\ \ge\ \sigma T,\qquad
           \sigma=\epsilon/4=1/1024 .
\end{equation}
\end{lemma}
\begin{proof}
Proceed from the bottom of the tree.
Fix such a node.
We bound the contribution
\(
   \sum_e\pi_e (\ln(1+\eta g_e)+L'_e)
\)
in~\eqref{eq:bonus} per probe.
\begin{itemize}
 \item The child succeeds:
 the contribution from log market gains
 \(\ln(1+\eta g_e)\)
 is nonnegative, and at least \(\epsilon\)
 if heavy, by Lemma~\ref{lem:gain}
 (the experts added for future pivots have gain zero).
 Child bonuses only help.
 \item The child fails:
 then the child must be internal by Lemma~\ref{lem:leaf}.
 Therefore
 \[
    \sigma r\le
         \sum_e \pi'_e L'_e
        \le C\sum_e \pi_e L'_e
 \]
 by induction and the choice of pivots. Thus this probe contributes at least
 \[
       \ln(1-\eta)+\frac{\sigma r}{C}
       \ge -2\eta+\frac{\sigma r}{C} \ge\epsilon ,
 \]
 paying for possibly bad market outcomes
 (again future-pivot experts have zero gain).
\end{itemize}
By the median rule, failure implies that for a majority
of probes either the probe was heavy or the child failed.
So altogether
\[
   \sum_e \pi_e L_e\ \ge\
      T\epsilon/2
          -3T(2c)-3\ln(2/c)\ \ge\ \sigma T
\]
by~\eqref{eq:bonus}, \(\ln(1-c)\ge-2c\),
and the choices~\eqref{eq:probepars}
(\(6c<\epsilon/8,\ 3\ln(2/c)<r\epsilon/8,\ T\ge r\)).
\end{proof}

At the root, \(\mu_e\le|E|\)
since the total wealth cannot increase,
so \(L_e\le\log(|E|+1)\)
(recall \(\log\) denotes base two).
As \(\sigma r_0>3\log(|E|+1)\),
there cannot be a failure: apply Lemma~\ref{lem:charge}
with assignments as in Lemma~\ref{lem:transfer-p}.
As this holds for any hypothesis consistent with all answers,
the returned prices meet all demands as required.

\subsection{Running the markets}\label{sec:implement-pr}

The number of probes is \(O(r_0r^{D-1})\). The expert choices give
\(|E|\) polynomial in \(N,R\), so \(r_0=O(\log(NR+1))\).
Since \(D\) is the least positive odd integer with \((5/4)^D>2N\)
and \(\log_{5/4}r<60\), the tree uses \(O(N^{60}\log(NR+1))\)
probes. The wealths also have polynomial bit length. Each probe adds
only constantly many binary places to a common dyadic denominator,
and all wealths are bounded by the root's initial total \(|E|\).

The remaining cost is choosing the proposal at each probe. We first
bound its bit length and then describe how to optimize over the
implicit expert pool in time polynomial in \(N,\log(R+1)\).

\paragraph{Prices at vertices.}
At each probe use an optimum vertex of the epigraph of~\(\Phi\).
Such vertices exist. Indeed, \(\Phi\) even tends to infinity
as \(\|y\|_\infty\) grows:
there are zero-threshold experts whose tuples
repeat a single slot with guess zero, one for each slot,
with strictly positive wealth.
Thus the optimal face is nonempty and bounded.

We claim this keeps all prices half-integral, with polynomial bit
length, through the tree. To see this, a vertex price vector
is determined by ties in each expert's option utilities, i.e.,
by equations equating utilities of options of the same expert
(including ties with the outside option).
Indeed, otherwise there is a line through the price vector
preserving all those ties; each \(U_e\), hence also \(\Phi\),
is affine near the point on this line, contradicting extremality.
Equations from equating two slot utilities have form
\[
      y_i = \pm y_j + B,\qquad B\in\ZZ,\quad |B|\le2R ,
\]
as the \(\tau\) is common. Within any connected component of these edges
on the \(N\) coordinates, express all its prices in terms of
one of them via a spanning tree.
If these equations in the component already determine the prices there,
some remaining edge (possibly a loop) determines twice
that one price as an integer of magnitude \(O(NR)\).
If not, a tie with the outside option must pin some price in the component,
to an integer or to a signed history price plus an integer
(in either case with the integer bounded by \(2R\)
in absolute value).
By induction down the tree we keep half-integrality
and prices of magnitude \(O(DNR)\).

\paragraph{Cells of experts.}
Group experts by all their data except their
demand guesses. Thus the number of groups is polynomial, and
each group corresponds to the integer points of a box
(one dimension per guess, at most five).
Observe that for given prices (proposal and history),
all comparisons among an expert's option utilities are
comparisons of single guesses or differences of two guesses
against constants. Thus we can subdivide the groups
by these comparisons, creating cells of experts with a common
set of maximizing options (by position on their lists).

More precisely, within each group partition
the range \([-R,R]\) for each guess into intervals of integers,
and similarly partition \([-2R,2R]\) for each pairwise difference.
Choosing one interval from each partition gives a cell of experts,
requiring all of the bounds simultaneously
(the cell may be empty). Start with each range as one interval.
Every time we choose proposal prices,
refine according to all of the option comparisons with these
prices, keeping ties separate.
Each partition now gets only a constant number of new cuts per probe,
and each group has only
a constant number of partitions.
Thus the number of cells remains polynomial.

The outcomes (including ties) of all option comparisons at
proposals so far are now common within cells.
So gains in response to feedback are also common:
we can identify all maximizing options at the proposal.
Wealths too are common, since splits, updates and recombination
preserve this.
Keep these per-expert data for each cell,
copying them upon refinement to the subcells.
This includes wealth data for any ongoing probes, which
we need to recombine wealths.
Compare Feodorov and Haslebacher's use of arrangements and volumes
for fixed points~\cite{FH};
here each expert has only a constant-dimensional guess,
regardless of \(N\).

To use these cells in choosing prices, we compute two kinds of sums:
the number of experts in a given cell, and the sum there of any
given guess coordinate.
In fact we can sum polynomials by eliminating variables one by one.
The range for a variable to eliminate runs from a maximum
of lower bounds to a minimum of upper bounds,
each bound an integer constant or a remaining variable plus
an integer constant.
So, on the remaining variables, we can subdivide according to which
bounds achieve these two extrema (choose first in case of ties),
and require lower at most upper.
All these conditions have the same form (interval constraints on
variables and differences). On each piece, summing over this range
again gives a polynomial by the formulas for sums of powers
(discrete antiderivatives, valid also with negative endpoints).
Starting with constant and linear polynomials for our computations,
the degree goes up by at most one per step, and bounds and coefficients
undergo only constantly many arithmetic operations:
both branching and depth are bounded by constants.
This gives the sums in polynomial bit complexity.

\paragraph{Computing the market minimizer.}
The label oracle is queried at prices returned by the children.
To choose the parent proposal, we construct a separate oracle from the
finite market data. It separates points from the epigraph of \(\Phi\),
which is defined by one inequality for every selection of one option
per expert: the height must dominate their weighted sum of utilities.
These inequalities have polynomial bit length because the history
prices are short, the weights share a short denominator, and \(\log|E|\)
is polynomial in the input size.

At rational test prices, compute an affine sum as above attaining
\(\Phi\).
To do so, temporarily refine on the option comparisons at
these prices. Choose a common maximizing option for the experts
in each cell and compute the coefficients from their weighted utilities
using counts and sums of guesses.
This gives separation: compare the proposed height to \(\Phi\)
and use the computed inequality if violated (the normal can be
scaled to sup norm one).
The bit time is polynomial, counting also the bits of the test point
as input.

Apply the classical bit-polynomial ellipsoid algorithm for
well-described polyhedra~\cite[Theorem~(6.4.9)]{GLS} to this separator.
All data passed to it have the bit bounds just proved. This is the
LP subroutine used to implement~\eqref{eq:market}.

Add the optimal height constraint
(still polynomial bits), and take
the lexicographically first optimal point,
a vertex (cf.~\cite[Lemma~(6.5.1)]{GLS}),
minimizing each coordinate successively.
This works in polynomial time as well:
the optimal face is bounded,
and each additional constraint value fixing a price coordinate
is a coordinate of the final vertex,
hence still short.

This gives the required deterministic implementation;
medians and labels too use polynomially many bits,
and the bounds hold regardless of labels.
This proves Theorem~\ref{thm:prices} and so completes
the proof of Theorem~\ref{thm:lp}.
